\documentclass[12pt]{article}
\usepackage[utf8]{inputenc}
\usepackage[T1]{fontenc}
\usepackage{hyperref}
\title{Yendang Language Resources for K-12 Curriculum Development}
\author{Compiled by Claude Sonnet 5.5 (claude-sonnet-5-5) for Andr\'as Kornai}
\date{}
\begin{document}
\maketitle

\section{Introduction}
The code \texttt{yen} is a retired ISO 639-3 code. It was split in 2012 into \texttt{ynq} (Yendang) and \texttt{yot} (Yotti). Glottolog notes that Yotti had wrongly been listed as a dialect of Yendang. Yendang is a Leko--Nimbari (Savanna) language of northeastern Nigeria. Wikipedia cites about 50,000 speakers (1987), and Global Recordings Network gives a population of 62,640 with no year. Resources for the language should be sought under the code \texttt{ynq}, not \texttt{yen}.

\section{Primary Educational Resources}
The searches found no primers, textbooks or school programmes for Yendang.

\section{Linguistic and Reference Materials}
\begin{itemize}
\item Glottolog entry for the retired code: \url{https://glottolog.org/resource/languoid/iso/yen}
\item Wikipedia, Yendang language: \url{https://en.wikipedia.org/wiki/Yendang_language}
\item IANA subtag record for ynq: \url{https://www.iana.org/assignments/lang-subtags-templates/ynq.txt}
\end{itemize}

\section{Religious and Community Resources}
No scripture or community organisation was found in the search results.

\section{Digital and Audio Resources}
Global Recordings Network lists Yendang audio material: \url{https://globalrecordings.net/en/language/ynq}

\section{Curriculum Development Context}
Little exists for this language. A curriculum would have to start from the \texttt{ynq} materials and fieldwork. The speaker figures are dated and should be checked.
\end{document}